\documentclass[11pt,a4paper]{letter}
\usepackage[utf8]{inputenc}
\usepackage[ngerman]{babel}
\usepackage{amsmath, amssymb}
\usepackage{csquotes}
\usepackage{geometry}
\geometry{margin=2.5cm}

\begin{document}
\begin{letter}

Dipl.-Ing. (FH) Michael Czybor \\ Hauptstraße 78b \\ 4222 Langenstein, AT \\ deutscher Staatsbürger

\vspace{0.5cm}

An die Fraktionen des Deutschen Bundestages

\begin{flushright}
	\today
\end{flushright}

\vspace{0.5cm}

Sehr geehrte Damen und Herren,

anliegend übermittle ich Ihnen eine Petition, die ich beim Deutschen Bundestag eingereicht habe. Ich bitte Sie, sich als Fraktion mit folgenden Fragen auseinanderzusetzen:

\begin{enumerate}
    \item Ist die Urknall-Hypothese noch haltbar, wenn ihre zentralen Säulen – dunkle Materie und dunkle Energie – nach über 50 Jahren intensivster Suche (CERN, LUX-ZEPLIN, XENONnT, ADMX, Euclid, Planck, James-Webb-Teleskop) nicht direkt nachgewiesen wurden?
    
    \item Welche wissenschaftlichen Alternativen (stationäre Kosmologien, Plasma-Kosmologie, fraktale Raumtheorien mit $D \approx 2,704$, Weber-Gravitation, de Broglie-Bohm-Quantentheorie, $\Omega$-Theorie) existieren jenseits des $\Lambda$CDM-Mainstreams, und warum werden diese Alternativen im deutschen Fördersystem – von der DFG über die Max-Planck-Gesellschaft bis zum BMBF – systematisch benachteiligt oder gar nicht erst zur Begutachtung zugelassen?
    
    \item Welche gesellschaftlichen, wirtschaftlichen und kulturellen Implikationen hat es, wenn ein ganzes Forschungs- und Weltbildparadigma (Urknall, Expansion, dunkle Entitäten, Hitzetod) auf prinzipiell nicht nachweisbaren Entitäten aufbaut – insbesondere im Hinblick auf die Parallele zwischen expandierendem Universum und expandierendem Wirtschaftssystem sowie zwischen Hitzetod und Klimakollaps?
\end{enumerate}

\noindent
\textbf{Zur beigefügten Arbeit}

Die Arbeit von Michael Czybor (\enquote{Die vereinheitlichte Theorie von Information, Raum und Gravitation} – IWT \& WDBT+) zeigt eine vollständige, konsistente, axiomatisch geschlossene Alternative ohne Urknall, ohne dunkle Materie, ohne dunkle Energie, ohne Inflation und ohne Singularitäten. Sie macht testbare Vorhersagen, die sich von denen des $\Lambda$CDM-Modells unterscheiden:

\begin{itemize}
    \item Frequenzabhängige Lichtablenkung ($\Delta\phi \propto \lambda^2$) – testbar mit EHT, VLBI, LISA
    \item Gravitationswellen-Dispersion ($v_{\text{phase}}(\omega) = c(1 + \alpha \omega^{-0,099})$) – testbar mit LISA, Einstein-Teleskop
    \item Maximale Masse kompakter Objekte $M_{\text{BEC}} \approx 3 \times 10^6 M_\odot$ – testbar durch Beobachtung von Galaxienzentren
    \item Neutrinomasse $m_\nu \approx 0,1\,\text{eV}/c^2$ – testbar mit KATRIN, JUNO, DUNE
    \item Fraktale Skalengesetze $\rho(r) \propto r^{-0,296}$, $v(r) \propto r^{0,352}$ – testbar mit EUCLID, Rubin-Observatorium (LSST)
\end{itemize}

Die strukturelle Parallele zwischen dem expandierenden Universum der Urknall-Kosmologie und dem expansiven Wachstumszwang der kapitalistischen Wirtschaft ist nicht zufällig: Beide Systeme basieren auf unsichtbaren, nicht nachweisbaren Entitäten (dunkle Energie / dunkle Materie vs. \enquote{unsichtbare Hand} des Marktes / Externalisierung von Kosten). Ebenso wenig zufällig ist die Parallele zwischen dem postulierten Hitzetod des Universums (Big Rip, Big Crunch, Big Freeze) und den selbstzerstörerischen Tendenzen der gegenwärtigen Umwelt- und Klimapolitik.

Eine stationäre, ewige Kosmologie – wie die in der beigefügten Arbeit dargestellte – bietet ein wissenschaftliches Fundament für Nachhaltigkeit, Kreislaufwirtschaft, Gleichgewicht und einen respektvollen Umgang mit der Natur. Sie entzieht dem apokalyptischen Narrativ der totalen Vernichtung (Big Rip, Hitzetod) die wissenschaftliche Grundlage und eröffnet eine neue Perspektive auf das Selbstverständnis des Menschen im Kosmos.

\noindent
\textbf{Meine konkreten Forderungen an Sie als Fraktion}

\begin{enumerate}
    \item Nehmen Sie die Petition formell zur Kenntnis und setzen Sie sich dafür ein, dass sie im zuständigen Ausschuss (Bildung, Forschung und Technikfolgenabschätzung) ordnungsgemäß behandelt wird – inklusive einer Anhörung von Sachverständigen, die beide Paradigmen ($\Lambda$CDM und Alternativen wie die beigefügte IWT/WDBT+) gleichermaßen repräsentieren.
    
    \item Bringen Sie das Thema \enquote{Alternative Kosmologien und die Krise des Urknall-Paradigmas} in Ihre fraktionsinternen Arbeitsgruppen für Wissenschaft, Forschung, Technologie und Grundsatzfragen ein.
    
    \item Setzen Sie sich für eine offene, interdisziplinäre öffentliche Debatte im Deutschen Bundestag ein, in der Sachverständige der etablierten Kosmologie ($\Lambda$CDM) und der alternativen Paradigmen (stationäre Kosmologie, Plasma-Kosmologie, Weber-Gravitation, fraktale Raummodelle, de Broglie-Bohm-Quantentheorie) gleichermaßen gehört werden – nicht als \enquote{Störfeuer}, sondern als legitime wissenschaftliche Positionen.
    
    \item Unterstützen Sie die Forderung nach einer unabhängigen, transparenten Kosten-Nutzen-Analyse der deutschen Beteiligung an Urknall-orientierten Großforschungsprojekten (CERN, LISA, Euclid, Planck, LIGO/Virgo/KAGRA, Einstein-Teleskop, SKA, James-Webb-Teleskop) im Vergleich zur Fusionsforschung und anderen anwendungsnahen Technologien – unter ausdrücklicher Einbeziehung der Opportunitätskosten und der wissenschaftlichen Alternativen.
    
    \item Fordern Sie von der Bundesregierung einen Bericht über das Verfahren der Mittelvergabe in der Grundlagenforschung (DFG, MPG, BMBF) mit speziellem Fokus auf die Behandlung von Außenseitermeinungen und alternativen Paradigmen – mit dem Ziel, methodischen Pluralismus als Qualitätskriterium zu verankern.
    
    \item Stellen Sie eine Kleine oder Große Anfrage an die Bundesregierung zu den in dieser Petition genannten Fragen – insbesondere zur Bewertung alternativer kosmologischer Modelle und zur Parallele zwischen Urknall-Kosmologie und kapitalistischem Wachstumszwang.
\end{enumerate}

\noindent
\textbf{Schlussbemerkung}

Es geht nicht um eine einfache Kürzung von Forschungsmitteln. Es geht um eine grundsätzliche Neuausrichtung: Weg von einem dogmatischen, expandierenden und apokalyptischen Weltbild – hin zu einem pluralistischen, nachhaltigen, zukunftsoffenen und demokratisch verfassten Wissenschaftsverständnis.

Die Entscheidung für oder gegen das Urknall-Paradigma ist keine rein physikalische Frage. Sie ist eine politische, kulturelle und weltanschauliche Grundsatzentscheidung mit erheblichen Auswirkungen auf das Selbstverständnis des Menschen, seinen Umgang mit der Natur, seine Wirtschaftsordnung und seine Zukunftsperspektive.

Die beigefügte Arbeit von Michael Czybor zeigt, dass eine Alternative existiert – konsistent, testbar, singularitätenfrei, ohne dunkle Entitäten, ohne Urknall, ohne Expansion. Sie verdient eine ernsthafte, öffentliche, politische und wissenschaftliche Debatte.

Ich bitte Sie um eine schriftliche Stellungnahme zu den aufgeworfenen Fragen und um eine Rückmeldung, ob und wie Sie das Anliegen in Ihrer Fraktion behandeln werden.

\vspace{1cm}

\begin{flushleft}
    \rule{6cm}{0.4pt}
    
    \vspace{0.2cm}
    
    Langenstein/AT, \today\\Michael Czybor
\end{flushleft}

\noindent
\textbf{Anlagen}

\begin{itemize}
    \item Petition an den Deutschen Bundestag (vollständiger Text)
    \item Czybor, Michael: \enquote{Die vereinheitlichte Theorie von Information, Raum und Gravitation} (IWT \& WDBT+), Juni 2026 – vollständiges Dokument
\end{itemize}

\end{letter}
\end{document}